\chapter{GENERAL PROJECT CONTEXT}

\section{Introduction}
The global context of this project is the growing need for modern data platforms capable of unifying and governing customer data at scale. Customer information is not limited to a single database: it is scattered across marketing systems, sales transactions, web events, product interaction logs, and customer support channels. Without a unified engineering layer, organizations face duplicated work, inconsistent metrics, low trust in dashboards, and slow experimentation cycles for analytics and machine learning.

This chapter presents the environment in which the project was conducted, the motivating problem, and the objectives and scope of CustomerDNA AI.

In addition to the technical challenge, customer data introduces important constraints on operational quality. A customer identifier may be missing or represented differently across sources. Dates may be encoded in multiple formats. Event streams may arrive with duplicates or late data. If these issues are not addressed early, they propagate through the pipeline and compromise the credibility of final analytics models. For this reason, this project treats \textbf{data quality} and \textbf{observability} as first-class requirements and integrates validation and monitoring into the platform, not as external afterthoughts.

\section{Host Organization and Project Environment}
This project is positioned in a professional environment focused on applied innovation in information systems, data engineering, and artificial intelligence. The host organization provides a practical context where data-driven solutions are expected to be robust, maintainable, and ready for operational execution.

In this context, the project emphasizes engineering discipline: containerized deployment, explicit configuration through environment variables, idempotent re-processing, and end-to-end observability. These aspects are crucial for turning a technical prototype into a reproducible platform that can be executed consistently across machines and environments.

\begin{figure}[htbp]
    \centering
    \includegraphics[width=0.65\textwidth,keepaspectratio]{\CompanyLogo}
    \caption{Company logo of Smart Automation Technologies (SAT)}
    \label{fig:company-logo}
\end{figure}

\subsection{Company Context and Digital Transformation Needs}
The host organization operates in a context where companies and public administrations increasingly rely on digital transformation to modernize their services and improve performance. In such environments, data becomes a strategic asset, but it is often fragmented across operational systems and external tools. Building value from data therefore requires a strong data engineering layer that can ingest, standardize, validate, and expose reliable datasets for analytics and decision-making.

From an organizational perspective, recurring needs include:
\begin{itemize}
    \item consolidating customer data from multiple business domains into a unified view
    \item enabling traceable, auditable, and reproducible data pipelines
    \item reducing the operational cost of reprocessing and troubleshooting data issues
    \item preparing structured, trusted datasets for advanced analytics and machine learning
\end{itemize}

In practice, digital transformation initiatives often begin with dashboards and reporting, but quickly evolve toward more advanced use cases such as churn prediction, customer lifetime value estimation, marketing campaign optimization, and personalization. These use cases require an engineering layer that is able to continuously produce curated datasets, not only a one-time extraction. For this reason, the platform built in this project is designed to be executed repeatedly, with clear separation between raw landing zones and refined analytical products.

In the context of the host organization and the project environment, the activities relevant to this work include support from requirements and architecture definition to implementation and operational follow-up. Typical delivery areas aligned with this project include:
\begin{itemize}
    \item \textbf{information systems and digital consulting}: requirements analysis, architecture design, and modernization roadmaps
    \item \textbf{data engineering and analytics}: ingestion pipelines, lakehouse organization, modeling, and BI-ready datasets
    \item \textbf{automation and operational improvement}: process automation and engineering of reliable operational data flows
    \item \textbf{artificial intelligence integration}: preparing data foundations for ML use cases and supporting experimentation workflows
    \item \textbf{capacity building}: technical training, workshops, and knowledge transfer for engineering teams
\end{itemize}

To ensure sustainable delivery, Smart Automation Technologies adopts an approach where engineering excellence is linked to ethics, transparency, and long-term partnership. This approach is particularly relevant for data projects, where trust is a strict requirement: data products must be traceable, transformations must be explainable, and results must be reproducible.

The relationship with stakeholders is built on proximity and continuous communication to better understand operational needs and anticipate future expectations. In the context of this project, this translates into aligning the platform with real execution constraints: multi-source integration, reprocessing capabilities, service readiness, and monitoring that supports fast diagnosis.

In addition, the company considers human capital as a key asset and invests in developing team skills through guidance, evaluation, and continuous learning. A data platform is not only a set of tools; it is also a shared methodology and a set of practices that teams must be able to maintain and evolve.

This philosophy is reflected through values that directly impact engineering quality and delivery outcomes:
\begin{itemize}
    \item \textbf{team spirit}: collaboration between data engineering, analytics, and operations to deliver consistent results
    \item \textbf{integrity}: rigorous validation, transparent reporting, and reliable pipeline traceability
    \item \textbf{commitment}: operational responsibility through repeatable execution and clear monitoring
    \item \textbf{excellence}: testing, documentation, and structured modeling to build trusted data products
    \item \textbf{innovation}: adoption of modern lakehouse practices to support scalable analytics and AI readiness
\end{itemize}

\begin{figure}[H]
    \centering
    \includegraphics[width=0.68\textwidth,keepaspectratio]{\CompanyValuesFigure}
    \caption{Core values of Smart Automation Technologies (SAT)}
    \label{fig:company-values}
\end{figure}

These values provide a natural framing for CustomerDNA AI. The platform is designed to produce reliable data products (integrity), to be executable and observable in operational conditions (commitment), to follow structured engineering patterns (excellence), to leverage modern scalable tooling (innovation), and to remain maintainable by teams through clear separation of concerns (team spirit).

\subsection{Data-Driven Decision-Making Challenges}
Data-driven decision-making requires more than data availability. Common difficulties observed in real-world contexts include inconsistent definitions of key metrics, missing lineage between raw sources and dashboards, limited scalability when data volume increases, and the absence of continuous validation and monitoring.

These challenges motivate a platform approach where ingestion, storage, transformation, and quality checks are treated as a coherent system. In this project, this coherence is ensured through a lakehouse design that separates responsibilities by layers (bronze, raw, staging, intermediate, analytics) and provides clear contracts between each stage.

In concrete terms, the project targets challenges such as:
\begin{itemize}
    \item \textbf{heterogeneous formats}: CSV files, Excel workbooks, and source-specific record structures
    \item \textbf{schema instability}: evolving columns and inconsistent naming conventions across sources
    \item \textbf{data consistency}: missing identifiers, duplicated events, invalid dates, and inconsistent types
    \item \textbf{reprocessing difficulty}: inability to reliably rebuild datasets after a bug or a logic change
    \item \textbf{limited observability}: lack of metrics about which datasets were loaded, when, and with what outcome
\end{itemize}

\subsection{AdOptimizer Global Project}
AdOptimizer is the broader project vision within which customer data, analytical intelligence, and future marketing optimization capabilities are organized. The present PFE does not implement the complete AdOptimizer product portfolio. It delivers the distributed data foundation required by those future products.

Within this global vision, CustomerDNA AI is responsible for transforming heterogeneous customer-related sources into governed, quality-aware, and queryable lakehouse products. This positioning preserves the connection to the AdOptimizer business objective while keeping the implemented academic contribution focused on Data Engineering.

\subsection{CustomerDNA AI and the Customer 360 Vision}
CustomerDNA AI is the distributed data engineering foundation for a broader Customer~360 and intelligent-marketing vision. It combines campaign-response, browsing-intent, and transaction signals so that future analytical and machine-learning consumers can work from governed and traceable tables.

The implemented contribution ends at the level of a deployed, observable, and quality-aware distributed lakehouse. Dynamic segmentation, churn prediction, LTV prediction, persona generation, and personalized recommendations are downstream possibilities rather than implemented PFE claims. This boundary keeps the project academically focused on Data Engineering while explaining its business value.

The Customer~360 perspective influences the design: source identity must be preserved, transformation logic must remain versioned, quality results must be auditable, and new sources must be onboarded through localized configuration rather than by rewriting unrelated layers.

\section{CustomerDNA AI Project Presentation}
The target project consists of building a distributed lakehouse platform that transforms raw customer-related datasets into structured, trusted, and queryable data products. The implemented solution is designed to be modular and reproducible, with each subsystem deployed using Docker and connected through well-defined interfaces.

The implemented Client~1 case study uses three customer-related datasets representing campaign response, digital browsing intent, and commercial transactions. This focused bundle remains heterogeneous while fitting the available distributed infrastructure and supporting a clear Customer~360 analytical narrative.

\subsection{CustomerDNA AI Platform Principle}
The core principle of the platform is to establish an end-to-end pipeline that is both scalable and operationally manageable. The platform follows a layered organization:
\begin{itemize}
    \item ingestion layer for publishing multi-source data into Kafka topics
    \item bronze layer for immutable storage of raw batches in HDFS
    \item raw lakehouse layer where Spark loads data into Iceberg tables managed by Hive Metastore
    \item transformation layers where dbt-Spark produces staging, intermediate, and analytics tables
    \item query layer where Trino enables SQL exploration and consumption
    \item validation and observability layers for quality checks and operational monitoring
\end{itemize}

This principle is validated through three complementary sources. Bank Marketing contributes customer profile, campaign-contact, response, and macro-economic attributes. Online Shoppers Purchasing Intention contributes session behavior, engagement, and purchase-intent signals. Online Retail~II contributes invoice, product, revenue, and repeat-purchase behavior. Each source follows the same Kafka, HDFS, Spark, Iceberg, dbt-Spark, quality, and orchestration contracts.

\begin{table}[htbp]
    \centering
    \small
    \setlength{\tabcolsep}{4pt}
    \renewcommand{\arraystretch}{1.06}
    \arrayrulecolor{ModernTableRule}
    \rowcolors{2}{ModernTableStripe}{white}
    \begin{tabularx}{0.96\textwidth}{|>{\raggedright\arraybackslash}p{0.25\textwidth}|>{\raggedright\arraybackslash}p{0.15\textwidth}|>{\raggedright\arraybackslash}p{0.23\textwidth}|>{\raggedright\arraybackslash}X|}
        \hline
        \rowcolor{ModernTableHeader}
        \color{white}\textbf{Dataset} &
        \color{white}\textbf{Source / format} &
        \color{white}\textbf{Iceberg raw table} &
        \color{white}\textbf{Main analytical purpose} \\
        \hline
        Bank Marketing & CSV & \path{raw_data.bank_marketing} & profiles, campaign response, and conversion \\
        \hline
        Online Shoppers Purchasing Intention & CSV & \path{raw_data.online_shoppers_intention} & session engagement and purchase intention \\
        \hline
        Online Retail II & Excel & \path{raw_data.online_retail_2} & transactions, sales value, and repeat purchase \\
        \hline
    \end{tabularx}
    \caption{Active Client 1 datasets used by CustomerDNA AI}
    \label{tab:case-study-datasets}
\arrayrulecolor{black}
\end{table}

The three sources have recognized academic provenance. Bank Marketing includes field and citation documentation in \path{bank-additional-names.txt}; Online Shoppers Purchasing Intention is a behavioral analytics dataset; and Online Retail~II is a widely used transactional retail dataset. Their documented origin and complementary grains strengthen the reproducibility and jury defensibility of the case study.

From a functional perspective, the end-to-end execution can be summarized as:
\begin{itemize}
    \item ingest source datasets and publish records into Kafka topics with consistent metadata
    \item consume Kafka topics, batch records, and land immutable JSONL files into an HDFS bronze zone
    \item trigger distributed Spark jobs that convert bronze batches into Iceberg raw tables
    \item model and curate data with dbt-Spark into staging, intermediate, and analytics layers
    \item query curated tables through Trino for exploration and downstream consumption
    \item validate each raw dataset through Great Expectations definitions, persist JSON results, generate Data Docs, and expose operational metrics
\end{itemize}

\subsection{Problem Statement}
The main problem addressed by this project is the absence of a unified and governed data engineering workflow for multi-source customer data. In practice, raw datasets are heterogeneous (CSV, Excel, event logs), arrive at different cadences, and may contain inconsistent types, missing identifiers, and invalid values.

Without a structured platform, teams typically rely on manual preprocessing or ad-hoc scripts, which creates recurring issues: lack of reproducibility, weak traceability, and limited scalability. This project therefore aims to deliver a system where ingestion, storage, transformation, and quality validation are standardized, automated, and observable.

More precisely, the problem can be decomposed into operational questions that a platform must answer reliably:
\begin{itemize}
    \item How can we ingest multiple datasets repeatedly while guaranteeing consistent results?
    \item How can we store raw data in a distributed system while keeping the ability to trace every transformation?
    \item How can we provide SQL access to both raw and curated data without exporting to a separate warehouse?
    \item How can we validate that the pipeline produced usable outputs and detect failures early?
    \item How can we monitor service readiness and pipeline execution in a way that supports debugging and reporting?
\end{itemize}

\subsection{Project Objectives}
The objectives of this project are both technical and operational:
\begin{itemize}
    \item design a distributed lakehouse architecture adapted to multi-source customer datasets
    \item implement an end-to-end pipeline from ingestion to analytics-ready tables
    \item ensure reproducibility through containerized deployment and explicit configuration
    \item provide data quality validation to improve trust in downstream analytics
    \item integrate orchestration and monitoring for production-oriented execution
\end{itemize}

In addition, the project aims to deliver a coherent repository structure and execution strategy that can be reused as a foundation for future extensions: onboarding additional datasets, adding new dbt models, integrating stronger data quality tests, and preparing feature stores or machine learning training pipelines.

\subsection{Project Scope and Boundaries}
The scope of the implemented platform is focused on data engineering and lakehouse construction. The project covers ingestion, distributed storage, distributed processing, transformation modeling, service readiness checks, and monitoring.

The platform is demonstrated using a defined set of customer-related datasets and is designed so that additional datasets and clients can be integrated by extending the configuration and pipeline scripts. Advanced machine learning model training and deployment are considered future extensions enabled by the resulting analytics-ready data products.

To keep the scope achievable and consistent with the objectives, the following boundaries are considered:
\begin{itemize}
    \item the project focuses on the engineering platform and curated analytical tables, not on delivering a final ML model in production
    \item security and access control are introduced as considerations and configuration practices, while full enterprise IAM integration is outside the scope
    \item the platform demonstrates distributed ingestion, storage, processing, and query execution; enterprise-scale benchmarking and advanced IAM hardening remain outside the implemented PFE scope
\end{itemize}

\subsection{Expected Outcomes}
At the end of the project, the expected outcomes include:
\begin{itemize}
    \item a working distributed pipeline that converts raw multi-source datasets into curated lakehouse tables
    \item a modular repository structure with separated layers and reproducible deployment procedures
    \item validated raw tables and tested transformation models, increasing trust in data products
    \item observable service health and pipeline execution metrics through Prometheus and Grafana
\end{itemize}

From a deliverables perspective, the expected outputs can be interpreted as:
\begin{itemize}
    \item a deployed centralized control-plane stack and three VM data-plane bundles for Kafka, HDFS, Spark, Hive Metastore, Trino, Airflow, Prometheus, and Grafana
    \item ingestion scripts that publish datasets to Kafka, with topic management and message confirmation
    \item bronze landing process that writes immutable batch files to HDFS and emits load events
    \item Spark job that loads bronze batches into Iceberg raw tables with verification steps
    \item dbt-Spark project producing multi-layer curated tables and associated tests
    \item dataset-specific Great Expectations validation definitions producing JSON artifacts and HTML Data Docs
    \item monitoring exporter exposing metrics for service readiness and pipeline execution history
\end{itemize}

\section{Project Management}

\subsection{Development Methodology and Project Planning}
The project methodology follows an engineering workflow organized around iterative delivery and validation. The main steps include requirements analysis, architecture design, technology selection, infrastructure implementation, incremental pipeline development, quality validation, and monitoring instrumentation.

Each major capability is implemented in a way that supports re-execution without manual cleanup: Kafka topics can be reset, the HDFS bronze zone can be re-initialized, raw tables can be rebuilt, and transformations can be re-run and tested. This approach ensures that the platform remains maintainable and verifiable throughout the report implementation and future experimentation.

In order to translate these steps into actionable engineering work, the project follows an incremental approach where every layer is validated before moving to the next:
\begin{itemize}
    \item validate the infrastructure layer by checking service readiness and connectivity
    \item validate ingestion by comparing produced message counts and consumed batch sizes
    \item validate storage by ensuring bronze files are created with stable naming and directory structure
    \item validate processing by comparing expected and actual row counts in raw tables
    \item validate transformations through dbt tests and curated table inspection
    \item validate quality and operations through monitoring metrics and dashboards
\end{itemize}

\begin{figure}[H]
    \centering
    \begin{adjustbox}{max width=\textwidth,center}
    \begin{tikzpicture}[
        node distance=10mm and 14mm,
        workflow/pill/.style={
            workflow/step,
            rounded corners=18pt,
            text width=3.0cm,
            minimum height=1.0cm,
            fill=white
        }
    ]
        \node (start) [workflow/pill] {Start};

        \node (tests) [workflow/step, fill=white, text width=3.35cm, below left=16mm and 34mm of start] {Run tests and\\validate execution};
        \node (analyze) [workflow/step, fill=white, text width=3.35cm, below=16mm of start] {Analyze needs\\and define\\objectives};
        \node (select) [workflow/step, fill=white, text width=3.35cm, below right=16mm and 34mm of start] {Select technologies\\and appropriate\\tools};

        \node (docs) [workflow/step, fill=white, text width=3.35cm, below=12mm of tests] {Produce\\documentation};

        \node (design) [workflow/step, fill=white, text width=3.35cm, below=22mm of analyze] {Design\\architecture};
        \node (impl) [workflow/step, fill=white, text width=3.35cm, below=22mm of select] {Implement\\infrastructure};

        \node (deploy) [workflow/step, fill=white, text width=3.1cm, below=22mm of docs] {Deploy};
        \node (opt) [workflow/step, fill=white, text width=3.1cm, below=10mm of deploy] {Optimize};
        \node (end) [workflow/pill, below=10mm of opt] {End of project};

        \begin{pgfonlayer}{background}
            \node (tasksbg) [
                draw=WorkflowOrange!55,
                line width=0.9pt,
                rounded corners=10pt,
                fill=WorkflowOrange!6,
                fit=(tests)(analyze)(select)(docs),
                inner sep=8mm
            ] {};
        \end{pgfonlayer}
        \node[anchor=north west, font=\normalsize, text=WorkflowOrange!80] at ([xshift=6mm,yshift=-4mm]tasksbg.north west) {TASKS};

        \node[draw=WorkflowBlue, fill=WorkflowBlue, circle, inner sep=1.6pt, text=white, font=\normalsize] at ([xshift=-1.55cm,yshift=3.5mm]tests.north) {1};
        \node[draw=WorkflowTeal, fill=WorkflowTeal, circle, inner sep=1.6pt, text=white, font=\normalsize] at ([xshift=-1.55cm,yshift=3.5mm]analyze.north) {2};
        \node[draw=WorkflowPurple, fill=WorkflowPurple, circle, inner sep=1.6pt, text=white, font=\normalsize] at ([xshift=-1.55cm,yshift=3.5mm]select.north) {3};

        \draw[workflow/arrow] (start) -- (analyze);
        \draw[workflow/arrow] (start.south) |- (select.north);
        \draw[workflow/arrow] (analyze.west) |- (tests.east);
        \draw[workflow/arrow] (tests) -- (docs);
        \draw[workflow/arrow] (docs) -- (deploy);
        \draw[workflow/arrow] (deploy) -- (opt);
        \draw[workflow/arrow] (opt) -- (end);
        \draw[workflow/arrow] (analyze) -- (design);
        \draw[workflow/arrow] (select) -- (impl);
        \draw[workflow/arrow] (impl.west) -- (design.east);
        \draw[workflow/arrow] (design.west) |- (deploy.east);
    \end{tikzpicture}
    \end{adjustbox}
    \caption{Global project workflow used to guide implementation}
    \label{fig:project-workflow}
\end{figure}

\clearpage
\subsection{Project Difficulties}
Implementing the distributed lakehouse platform introduced four main engineering challenges:

\begin{enumerate}
    \item \textbf{Technology compatibility.} Coordinating versions, configurations, and network connectivity across Kafka, HDFS, Spark, Iceberg, Trino, Airflow, and the remaining services required continuous verification.
    \item \textbf{End-to-end integration.} Connecting ingestion, storage, processing, transformation, quality, and monitoring layers required clear interfaces and systematic validation at every stage.
    \item \textbf{Limited computing resources.} The three active datasets and explicit service limits were selected to demonstrate real distribution across three resource-constrained virtual machines.
    \item \textbf{Fragmented documentation.} Differences between documented deployment contexts complicated configuration and troubleshooting, making precise project documentation essential for reproducibility.
\end{enumerate}

\section{Conclusion}
This chapter established the AdOptimizer context, the Customer~360 objective, and the implemented boundaries of CustomerDNA AI. Chapter~II presents the theoretical foundations, Chapter~III formalizes the requirements, and Chapter~IV justifies the selected technology stack before design, implementation, and evaluation.